\documentclass[12pt,a4paper]{article}
\usepackage[utf8]{inputenc}
\usepackage[vietnamese]{babel}
\usepackage{amsmath,amssymb,amsthm}
\usepackage{geometry}
\usepackage{enumitem}
\usepackage[most]{tcolorbox}
\usepackage{hyperref}
\usepackage{fancyhdr}
\usepackage{tikz}
\usepackage{eso-pic}
\usepackage{xcolor}
\geometry{a4paper, margin=2cm, bottom=2.5cm}
\hypersetup{colorlinks=true, linkcolor=blue!70!black}
\setlist[itemize]{topsep=4pt, itemsep=2pt}
\setlist[enumerate]{topsep=4pt, itemsep=2pt}
\AddToShipoutPictureFG{%
  \AtPageCenter{%
    \begin{tikzpicture}[remember picture, overlay]
      \node[rotate=45, scale=1, opacity=0.06]{\fontsize{60}{72}\selectfont\bfseries\sffamily Đặng Tấn Quí};
    \end{tikzpicture}%
  }%
}
\pagestyle{fancy}
\fancyhf{}
\fancyhead[L]{\small\textit{TOÁN CHO KHOA HỌC XÃ HỘI}}
\fancyhead[R]{\small\textit{Đặng Tấn Quí}}
\fancyfoot[C]{\thepage}
\fancyfoot[L]{\footnotesize\textcopyright\ Đặng Tấn Quí}
\fancyfoot[R]{\footnotesize SĐT: 0377\,122\,210}
\renewcommand{\headrulewidth}{0.4pt}
\renewcommand{\footrulewidth}{0.4pt}
\fancypagestyle{plain}{\fancyhf{}\fancyfoot[C]{\thepage}\fancyfoot[L]{\footnotesize\textcopyright\ Đặng Tấn Quí}\fancyfoot[R]{\footnotesize SĐT: 0377\,122\,210}\renewcommand{\headrulewidth}{0pt}\renewcommand{\footrulewidth}{0.4pt}}
\newtcolorbox{dinhnghia}[1][]{enhanced, breakable, colback=blue!3!white, colframe=blue!60!black, fonttitle=\bfseries, title={Định nghĩa}, #1}
\newtcolorbox{dinhly}[1][]{enhanced, breakable, colback=green!3!white, colframe=green!50!black, fonttitle=\bfseries, title={Định lý}, #1}
\newtcolorbox{tinhchat}[1][]{enhanced, breakable, colback=yellow!5!white, colframe=orange!50!black, fonttitle=\bfseries, title={Tính chất}, #1}
\newtcolorbox{phuongphap}[1][]{enhanced, breakable, colback=gray!5!white, colframe=gray!50!black, fonttitle=\bfseries, title={Phương pháp}, #1}
\newtcolorbox{luuy}[1][]{enhanced, breakable, colback=red!3!white, colframe=red!50!black, fonttitle=\bfseries, title={Lưu ý}, #1}
\newtcolorbox{congthuc}[1][]{enhanced, breakable, colback=cyan!3!white, colframe=cyan!50!black, fonttitle=\bfseries, title={Công thức}, #1}
\newtcolorbox{vidu}[1][]{enhanced, breakable, colback=violet!3!white, colframe=violet!50!black, fonttitle=\bfseries, title={Ví dụ}, #1}

\begin{document}
\thispagestyle{empty}
\begin{tikzpicture}[remember picture, overlay]
  \fill[blue!80!black] (current page.south west) rectangle (current page.north east);
  \node[anchor=west, white, font=\fontsize{26}{32}\bfseries\selectfont]
    at ([xshift=3cm, yshift=-7.5cm]current page.north west) {TOÁN CHO KHOA HỌC XÃ HỘI};
  \node[anchor=west, white, font=\fontsize{18}{24}\selectfont]
    at ([xshift=3cm, yshift=-9cm]current page.north west) {Xác suất, thống kê, Markov, Bayes, hồi quy ứng dụng};
  \node[anchor=west, white, font=\Large\bfseries] at ([xshift=3cm, yshift=-13cm]current page.north west) {Biên soạn:};
  \node[anchor=west, yellow!80!white, font=\fontsize{22}{28}\bfseries\selectfont] at ([xshift=3cm, yshift=-14.5cm]current page.north west) {ĐẶNG TẤN QUÍ};
  \node[anchor=south east, white, font=\Large\bfseries] at ([xshift=-2cm, yshift=3cm]current page.south east) {\the\year};
\end{tikzpicture}
\newpage
\tableofcontents
\newpage
%% ================================================================
\section{Toán cơ bản ứng dụng}

\begin{congthuc}[title={Hàm tuyến tính và bậc hai}]
  $y = ax + b$: mô hình tuyến tính đơn giản (xu hướng, tỉ lệ).

  $y = ax^2 + bx + c$: đỉnh tại $x = -b/(2a)$ (tối ưu sản lượng, chi phí).
\end{congthuc}

\begin{congthuc}[title={Hàm mũ và logarit}]
  Tăng trưởng: $P(t) = P_0 e^{rt}$. Tăng trưởng logistic: $P(t) = \dfrac{K}{1 + Ae^{-rt}}$.

  Bán rã: $N(t) = N_0 e^{-\lambda t}$, $t_{1/2} = \ln 2/\lambda$.
\end{congthuc}

\begin{phuongphap}[title={Phần trăm và chỉ số}]
  Thay đổi phần trăm: $\Delta\% = \dfrac{x_{\text{new}} - x_{\text{old}}}{x_{\text{old}}} \times 100\%$.

  Chỉ số giá: Laspeyres, Paasche, Fisher.
\end{phuongphap}

%% ================================================================
\section{Ma trận và ứng dụng}

\begin{dinhnghia}[title={Xích Markov rời rạc}]
  Ma trận chuyển $\mathbf{P}$: $p_{ij} = \Pr(X_{n+1} = j \mid X_n = i)$.

  Phân phối sau $n$ bước: $\boldsymbol{\pi}^{(n)} = \boldsymbol{\pi}^{(0)}\mathbf{P}^n$.

  Phân phối dừng: $\boldsymbol{\pi} = \boldsymbol{\pi}\mathbf{P}$, $\sum \pi_i = 1$.
\end{dinhnghia}

\begin{vidu}[title={Ứng dụng}]
  Di cư dân số giữa các vùng, thay đổi thị phần, mô hình bầu cử.
\end{vidu}

%% ================================================================
\section{Xác suất ứng dụng}

\begin{dinhnghia}[title={Xác suất cơ bản}]
  $P(A) \in [0,1]$, $P(\Omega) = 1$. $P(A \cup B) = P(A) + P(B) - P(A \cap B)$.

  Xác suất có điều kiện: $P(A|B) = P(A \cap B)/P(B)$.
\end{dinhnghia}

\begin{dinhly}[title={Bayes}]
  $P(H|D) = \dfrac{P(D|H)P(H)}{P(D)}$. Ứng dụng: chẩn đoán, phân loại, cập nhật niềm tin.
\end{dinhly}

\begin{dinhnghia}[title={Phân phối phổ biến}]
  Binomial (bầu cử, khảo sát), Poisson (tội phạm, tai nạn), Normal (điểm thi, thu nhập biến đổi).
\end{dinhnghia}

%% ================================================================
\section{Thống kê mô tả và suy diễn}

\begin{phuongphap}[title={Mô tả dữ liệu}]
  Trung bình, trung vị, mode. Phương sai, độ lệch chuẩn, IQR. Biểu đồ: histogram, box plot, scatter.
\end{phuongphap}

\begin{phuongphap}[title={Khoảng tin cậy}]
  Trung bình: $\bar{x} \pm z_{\alpha/2}\dfrac{s}{\sqrt{n}}$ (mẫu lớn).

  Tỉ lệ: $\hat{p} \pm z_{\alpha/2}\sqrt{\dfrac{\hat{p}(1-\hat{p})}{n}}$.
\end{phuongphap}

\begin{phuongphap}[title={Kiểm định giả thuyết}]
  $H_0$ vs $H_1$, p-value, mức ý nghĩa $\alpha$. Chi-squared test cho bảng chéo (independence, goodness-of-fit).
\end{phuongphap}

%% ================================================================
\section{Hồi quy ứng dụng cho KHXH}

\begin{dinhnghia}[title={Hồi quy tuyến tính}]
  $Y = \beta_0 + \beta_1 X + \varepsilon$. Diễn giải: $\beta_1$ --- thay đổi $Y$ khi $X$ tăng 1 đơn vị. $R^2$: tỉ lệ giải thích.
\end{dinhnghia}

\begin{dinhnghia}[title={Hồi quy logistic}]
  Biến nhị phân: $\ln\dfrac{p}{1-p} = \beta_0 + \beta_1 X$. $e^{\beta_1}$: odds ratio.
\end{dinhnghia}

\begin{luuy}
  Tương quan không phải nhân quả. Biến gây nhiễu (confounding), bias chọn mẫu. Thiết kế thí nghiệm ngẫu nhiên (RCT) mới cho phép kết luận nhân quả.
\end{luuy}

\end{document}